\usepackage[letterpaper,margin=0.95in,headheight=15pt]{geometry}
\usepackage[T1]{fontenc}
\usepackage[utf8]{inputenc}
\usepackage{amsmath,amsthm,mathtools}
\usepackage{newpxtext,newpxmath}
\usepackage{microtype}
\usepackage{graphicx,booktabs,longtable,array,tabularx}
\usepackage{ragged2e}
\newcolumntype{P}[1]{>{\RaggedRight\arraybackslash}p{#1}}
\renewcommand{\tabularxcolumn}[1]{>{\RaggedRight\arraybackslash}p{#1}}
\usepackage{xcolor}
\definecolor{ink}{HTML}{143A45}
\definecolor{soft}{HTML}{F1F5F5}
\usepackage{enumitem}
\setlist{topsep=4pt,itemsep=2pt,parsep=0pt}
\usepackage{listings}
\lstset{basicstyle=\ttfamily\footnotesize,breaklines=true,columns=fullflexible,keepspaces=true,showstringspaces=false,frame=single,rulecolor=\color{black!20},xleftmargin=5pt,xrightmargin=5pt,aboveskip=9pt,belowskip=9pt}
\usepackage{tikz}
\usetikzlibrary{arrows.meta,positioning,calc,fit}
\usepackage[most]{tcolorbox}
\newtcolorbox{scopebox}{breakable,colback=soft,colframe=ink!45,boxrule=0.5pt,arc=0pt,left=9pt,right=9pt,top=7pt,bottom=7pt}
\usepackage{titlesec}
\titleformat{\section}{\Large\bfseries\color{ink}}{\thesection}{0.65em}{}
\titleformat{\subsection}{\large\bfseries}{\thesubsection}{0.6em}{}
\titlespacing*{\section}{0pt}{20pt}{9pt}
\titlespacing*{\subsection}{0pt}{12pt}{6pt}
\usepackage{fancyhdr}
\pagestyle{fancy}\fancyhf{}\fancyhead[L]{\small\textsc{\runningtitle}}\fancyhead[R]{\small Research edition 1.0}\fancyfoot[C]{\thepage}
\renewcommand{\headrulewidth}{0.3pt}
\usepackage[colorlinks=true,linkcolor=ink,citecolor=ink,urlcolor=ink,hypertexnames=false]{hyperref}
\usepackage[nameinlink,noabbrev]{cleveref}
\newtheorem{theorem}{Theorem}[section]
\newtheorem{proposition}[theorem]{Proposition}
\newtheorem{lemma}[theorem]{Lemma}
\newtheorem{corollary}[theorem]{Corollary}
\theoremstyle{definition}
\newtheorem{definition}[theorem]{Definition}
\newtheorem{example}[theorem]{Example}
\newtheorem{counterexample}[theorem]{Counterexample}
\theoremstyle{remark}\newtheorem{remark}[theorem]{Remark}
\newcommand{\R}{\mathbb R}
\newcommand{\N}{\mathbb N}
\newcommand{\im}{\operatorname{im}}
\newcommand{\rank}{\operatorname{rank}}
\newcommand{\id}{\operatorname{id}}
\newcommand{\diag}{\operatorname{diag}}
\newcommand{\norm}[1]{\left\lVert#1\right\rVert}
\newcommand{\code}[1]{\texttt{#1}}
\newcommand{\status}[1]{\textbf{#1}}
\setlength{\parindent}{0pt}
\setlength{\parskip}{5.5pt}
\renewcommand{\arraystretch}{1.13}
\emergencystretch=2em
\raggedbottom
\setcounter{tocdepth}{1}
\usepackage{tocloft}
\setlength{\cftbeforesecskip}{1.5pt}
\renewcommand{\cftsecfont}{\small}
\renewcommand{\cftsecpagefont}{\small}
\usepackage{needspace,etoolbox}
\pretocmd{\section}{\Needspace{7\baselineskip}}{}{}
\pretocmd{\subsection}{\Needspace{5\baselineskip}}{}{}
\ifdefined\pdfinfoomitdate\pdfinfoomitdate=1\fi
\ifdefined\pdftrailerid\pdftrailerid{}\fi
\ifdefined\pdfsuppressptexinfo\pdfsuppressptexinfo=15\fi
\newcommand{\editiontitle}[3]{%
\begin{titlepage}\thispagestyle{empty}
\vspace*{0.55in}
{\small\sffamily\bfseries\color{ink} KOMOREBI\quad /\quad LANGUAGE AND COMPUTATION}\par
\vspace{0.8in}
{\fontsize{35}{39}\selectfont\bfseries #1\par}
\vspace{0.28in}
{\Large #2\par}
\vspace{0.5in}
\rule{\textwidth}{0.6pt}\par
\vspace{0.2in}
{\large Komorebi P.}\par
Consolidated research edition 1.0\par
6 October 2026\par
\vspace{0.45in}
\begin{scopebox}#3\end{scopebox}
\vfill
{\small The manuscript, reference implementation, preserved source archives, and validation records form one edition. Publication finalization does not convert finite evidence into an unrestricted theorem or an unrun experiment into a result.}\par
\end{titlepage}}
